\begin{abstract}
In this study, we present an analysis of stable inorganic materials suitable for use as high-voltage protective barrier coatings near power lines. Criteria for selection included energy above hull between 0.0 and 0.05 eV/atom, band gap between 3.0 and 20.0 eV, and density between 3.0 and 10.0 g/cm$^{3}$, with a focus on excluding materials containing toxic elements such as Pb, Cd, and Hg. A thorough examination of retrieved materials data is complemented by visual analysis using scatter plots and histograms.
\end{abstract}

\section{Introduction}
The need for effective high-voltage protective barrier coatings in power line environments is critical due to their role in enhancing safety and reliability. Such coatings must exhibit a balance of stability, characterized by low energy above hull values, and electrical insulation, suggested by large band gaps, along with manageable densities that ensure mechanical integrity without excessive mass.

\section{Methods}
For this investigation, we leveraged data retrieval techniques to access a database of inorganic materials, filtering based on the criteria of stability, electronic properties, and non-toxicity. The selection criteria included energy above hull values from 0.0 to 0.05 eV/atom, band gaps from 3.0 to 20.0 eV, and densities ranging from 3.0 to 10.0 g/cm$^{3}$. Materials containing Pb, Cd, and Hg were excluded to ensure environmental and health safety. Data analysis was supported by visual tools to uncover trends and suitability of the materials.

\section{Results}
A total of 100 materials were identified that meet the specified criteria. The distributions of their properties were visualized through several analytic plots. Three scatter plots were constructed to explore relationships: band gap versus energy above hull, density versus band gap, and density versus energy above hull. Additionally, histograms displayed the distribution of band gaps and energies above hull. Most materials exhibited wide band gaps essential for insulation, and their energy above hull values were predominantly low, indicating high stability.

\begin{figure}[H]
\centering
\includegraphics[width=0.8\textwidth]{scatter_plot_of_band_gap_vs_energy_above_hull.png}
\caption{Scatter plot of band gap vs energy above hull, illustrating the relationship between electronic properties and material stability.}
\end{figure}

\begin{figure}[H]
\centering
\includegraphics[width=0.8\textwidth]{scatter_plot_of_density_vs_band_gap.png}
\caption{Scatter plot of density vs band gap, highlighting tradeoffs between mechanical and electronic properties.}
\end{figure}

\begin{figure}[H]
\centering
\includegraphics[width=0.8\textwidth]{scatter_plot_of_density_vs_energy_above_hull.png}
\caption{Scatter plot of density vs energy above hull, emphasizing correlations between mechanical stability and material density.}
\end{figure}

\begin{figure}[H]
\centering
\includegraphics[width=0.8\textwidth]{histogram_of_band_gap.png}
\caption{Histogram of band gap distribution among selected materials.}
\end{figure}

\begin{figure}[H]
\centering
\includegraphics[width=0.8\textwidth]{histogram_of_energy_above_hull.png}
\caption{Histogram illustrating energy above hull distribution, indicating material stability.}
\end{figure}

\section{Discussion}
The analysis has identified a diverse range of materials exhibiting properties favorable for high-voltage environment applications. The predominance of wide band gaps across the sample supports robust electrical insulation performance, while low energy above hull positions these materials within a stable regime suitable for deployment. The distribution of densities suggests that the mechanical implications of these materials in potential applications should be manageable.

\section{Conclusion}
This work highlights a subset of inorganic materials that meet critical thresholds for use as high-voltage protective barrier coatings, fulfilling requirements for stability and electronic insulation while avoiding toxic elements. Future studies may focus on integrating thermal and mechanical testing for a holistic assessment of material performance in real-world conditions.
